\begin{table}[tb]
\centering\footnotesize\setlength{\tabcolsep}{3pt}
\caption{保存入力の全６条件．各３予定生成．公開一致は既知使用量の比較可能対だけの一致数である．}
\begin{tabular}{lrrrr}
入力 & 試行 & 既知 & 完全非空 & 公開一致 \\ \hline
A & 3/3 & 3/3 & 3/3 & 3/3 \\
B & 3/3 & 3/3 & 3/3 & 3/3 \\
C & 3/3 & 3/3 & 3/3 & 3/3 \\
D & 3/3 & 3/3 & 3/3 & 3/3 \\
N & 3/3 & 3/3 & 3/3 & 3/3 \\
X & 3/3 & 3/3 & 3/3 & 3/3 \\
\hline\end{tabular}
\end{table}
この保存文脈では元のID一対で公開出力が異なり、同値入力の対照は一致した．
A/Bの公開出力一致は0/9対であった．
C/Dの公開出力一致は9/9対であった．
A/Nの公開出力一致は9/9対であった．
全18予定枠中18生成を試行し，既知422,475論理token，使用量不明0生成を記録した．所有推論器の稼働は4,386.052秒であった．全15組の比較と返却message全体の不透明hash照合を補足記録に保持し，内部推論の意味内容は読まない．一致する出力が正しいことや，継続後の品質が保たれることは判定していない．
